\section{The anchor construction}\label{sec:anchors}

The common-bag pseudometric is one part of the lower bound. We now
construct a second cut pseudometric $D_{\mathrm{anc}}$ by launching
particle pairs at bag occurrences and retaining only boundedly many
pairs along each descending path. An \emph{anchor} is a bag occurrence
together with the reserved pair launched there. Distinct occurrences of
the same vertex set give distinct possible anchors. We assume $|V|\geq2$
and use nonempty bags of size at most $k$; for a one-vertex graph both
pseudometrics are zero.

At a launch bag $A$, the update produces $f_A+g_A$ and $f_A-g_A$,
where $f_A$ is the average of the selected particles, and $g_A$
vanishes on $A$ and is at least $d(\,\cdot\,,A)$.
If $A$ lies on the bag path between two tested vertices and its pair
survives to the terminal bag, the launch-uncertainty test forces
separation. An older anchor may instead lie above the least common ancestor
of the two endpoint bags. Its launch bag is then unavailable for this test. If its
reservation reaches the terminal bag, we can use a later testing bag
on the path and seek a height gap there.

We obtain these later height gaps from two independent parts of $g_A$:
a random multiple of $h_A=d(\,\cdot\,,A)$, and a smaller localized
noise function. For a tested endpoint $x$ and a vertex $s$ of a testing
bag, the first coefficient varies their height difference when
$h_A(x)-h_A(s)$ is appreciable. A local coefficient varies the height
at $x$ without changing it at $s$ outside the bump's support. The next
lemma supplies both variations with uniform Lipschitz control. We then
give the deterministic retention rule. The testing bags used in
Section~\ref{sec:lower} will be chosen from the metric and that rule,
independently of the launch variables.

\subsection{Two independent sources of launch noise}

\begin{lemma}[Ordinary launch noise]\label{lem:ordinary-noise}
Let $A\subseteq V$ be nonempty, put $h(z)=d(z,A)$, and let
$0<\eta<1/4$. One can sample a random function
\[
  g=\xi h+v,
\]
using randomness independent of any previously sampled data. The
function $v$ is constructed by first sampling a finite indexed family
of nonnegative functions $(\varphi_i)_{i\in\mathcal I}$ on $V$,
called bumps, and then setting
\[
 v=\sum_{i\in\mathcal I}U_i\varphi_i.
\]
Conditional on the bump family, the coefficients
$U_i$ are independent and uniform on $[0,1]$. The following properties
hold:
\begin{enumerate}
\item $\xi$ is uniform on $[1,2]$ and independent of the bump family
      and all its coefficients;
      $v$ is $1$-Lipschitz and satisfies $0\leq v\leq\eta h/2$.
      Consequently $g|_A=0$, $\operatorname{Lip}(g)\leq3$, and
      \[
        h\leq g\leq \tfrac52h\leq\tfrac52\Delta.
      \]
\item Set $c=\alpha/8$ and $\rho=\pi/2$. For each fixed $x$ with
      $D_0=h(x)>0$, there is an event of probability at least $\rho$,
      determined by the bump family, on which an index $i(x)\in\mathcal I$,
      also determined by that family, selects a bump $\varphi_x=\varphi_{i(x)}$ satisfying
      \[
        \varphi_x(x)\geq c\eta D_0,
        \qquad
        \varphi_x(z)=0\quad\text{if }d(x,z)>4\eta D_0.
      \]
      Write $U_x=U_{i(x)}$. Conditional on the bump family, $U_x$ is uniform and independent
      of $\xi$, of the other bump coefficients, and of all previously
      sampled data. In particular,
      $v=v^{\setminus x}+U_x\varphi_x$, where the remaining coefficients
      defining $v^{\setminus x}$ are independent of $U_x$.
\end{enumerate}
Different vertices may select the same bump and coefficient. The
constants in these assertions are independent of $\eta$, $A$, the
number of vertices, and the ratio of the edge lengths. Only finitely many
random partitions and coefficients are required.
\end{lemma}

\begin{proof}
We divide the positive values of $h$ into randomly shifted dyadic
bands and put nonnegative bumps in partition parts within each band.
Conditional on their locations, we give the bumps independent
coefficients.  Disjoint supports keep their sum Lipschitz.  A vertex well
inside its band and its partition part then has one coefficient whose
amplitude is a fixed fraction of $\eta h(x)$.

If $h$ vanishes everywhere, take $v=0$ and sample $\xi$ as stated.
Otherwise let $h_{\min}$ and $h_{\max}$ be the smallest positive and the
largest values of $h$. Sample $\lambda$ uniformly on $[0,1)$, and put
\[
  I=\{\lfloor\log_2h_{\min}\rfloor-1,\ldots,
            \lfloor\log_2h_{\max}\rfloor\},
  \qquad m=\lfloor\log_2\eta\rfloor.
\]
For $j\in I$, define the random band and the fixed scale
\[
  W_j=\{z:2^{j+\lambda}\leq h(z)<2^{j+1+\lambda}\},
  \qquad r_j=2^{j+m}.
\]
The bands partition $\{h>0\}$. Independently of $\lambda$, sample a
partition $\mathcal P_j$ at scale $r_j$ from
Theorem~\ref{thm:padded}, for each $j\in I$. Thus the partition scales
are fixed powers of two, not random functions of $\lambda$.
Consequently, conditioning on the band offset leaves each sampled
partition with its original padding guarantee.

For each nonempty $E=W_j\cap C$, $C\in\mathcal P_j$, set
\[
  \varphi_{j,C}(z)=
  \begin{cases}
    \min\{r_j/2,d(z,V\setminus E)\},&z\in E,\\
    0,&z\notin E.
  \end{cases}
\]
Conditional on these locations, give the bumps independent uniform
coefficients $U_{j,C}\in[0,1]$, and define
\[
  v(z)=\sum_{j,C}U_{j,C}\varphi_{j,C}(z).
\]
This is a finite measurable construction: there are finitely many bands
and partitions of the finite set $V$, and band memberships are measurable
functions of $\lambda$. If desired, the coefficients can be sampled in
advance for every possible pair $(j,C)$ with $C\subseteq V$.

The nonempty sets $E$ are disjoint, and each bump is nonnegative and
$1$-Lipschitz on all of $V$. Lemma~\ref{lem:disjoint-bumps} therefore
shows that $v$ is $1$-Lipschitz, including at points where $h=0$. For $z\in W_j$,
\[
  0\leq v(z)\leq r_j/2
       \leq\eta 2^j/2\leq\eta h(z)/2.
\]
Sampling $\xi$ independently gives the first assertion.

Fix $x$ with $D_0=h(x)>0$, and let $j$ be its band index. The event
\[
  \log_2D_0-(j+\lambda)\in[1/4,3/4]
  \tag{*}
\]
has probability $1/2$. On this event, the value $D_0$ is at distance
greater than $D_0/8$ from both numerical endpoints of its band. Since
$h$ is $1$-Lipschitz, the ball $B_d(x,D_0/8)$ is contained in $W_j$.
Also $r_j\leq\eta D_0$, so $\alpha r_j<D_0/8$.
Conditional on $\lambda$, the probability that
$B_d(x,\alpha r_j)$ lies in its part $C\in\mathcal P_j$ is at least
$\pi$. On this padding event together with $(*)$,
\[
  \varphi_{j,C}(x)\geq\alpha r_j.
\]
Because $2^m>\eta/2$ and $D_0<2^{j+2}$, we have
$r_j>\eta D_0/8$. Thus the last display is at least
$c\eta D_0$ with $c=\alpha/8$. The joint event has probability at
least $\rho=\pi/2$ and depends only on the locations.
Finally, any point in the support of this bump is in the same partition
part as $x$, and hence is at distance at most
$r_j\leq\eta D_0$ from $x$. This implies the stated support bound.
Let $E_x$ be the event that the sampled family contains a bump with
the two displayed bounds in the second assertion, and select the first such
bump in a fixed ordering. This event and the selected index $i(x)$
are determined by the family. The band and padding event above implies
$E_x$, so $\Pr(E_x)\geq\rho$. The independent coefficients are
sampled after the family, which gives the stated conditional law of
$U_x=U_{i(x)}$. Finally, $\xi$ is sampled independently of both the
family and all coefficients.
\end{proof}

\subsection{Retaining anchors across distance drops}

For a vertex $p$ of the current bag, its distances to earlier launch
bags decrease with launch time, by the separator property. Our rule
protects the newer entry when two consecutive distances cross a
threshold determined by a distance $d(p,q)$ inside the current bag.
Two thresholds for each ordered pair $(p,q)$ protect only boundedly
many entries. In the lower-bound argument, a distance visible in the
bag will force this rule to retain the immediate successor of an old
anchor. For now let $J>0$ and $0<\eta<1/4$ be parameters; the memory
bound will be independent of their values.


Choose
\[
 n_0=2k^2+2.
\]
The cache is ordered by launch time, from oldest to youngest.  Launch an
anchor at the root and retain it permanently.  At every other bag $B$,
propose a new anchor launched
at $B$.  If the resulting list has at most $n_0$ entries, retain it.
Otherwise it has $n_0+1$ entries, and exactly one non-root entry is evicted;
the proposed anchor itself is an allowed eviction.  Copies of a parent's
cache and particle list are maintained separately for its children.

Here is the eviction rule, with a parameter $J>0$ to be fixed below.
For every ordered pair of distinct vertices $p,q\in B$, consider the two
thresholds
\begin{equation}\label{eq:cache-thresholds}
 \frac{d(p,q)}{100},\qquad Jd(p,q).
\end{equation}
For each threshold $\tau$, protect the entry immediately following a drop
from a distance strictly greater than $\tau$ to a distance at most $\tau$
in the list
\[
 d(p,A_1),d(p,A_2),\ldots,d(p,A_m),
\]
where the $A_i$ are the anchors, including the proposal, in age order.
If there is no such drop, that threshold protects no entry.  Evict the
oldest unprotected non-root entry.  All tie-breaking choices made here
are deterministic functions of the metric and decomposition.
In particular, the launch and release times of every reservation are fixed
before any perturbation is sampled.  The identities of the two particles
used at a launch may depend on earlier randomness, but the first-eligible-pair
rule of Proposition~\ref{prop:particle-realization} selects them before
sampling that launch's fresh variables.

\begin{lemma}[The cache is well defined]\label{lem:cache-defined}
The displayed distance list is nonincreasing, and an eligible eviction
always exists.  The cache has at most $n_0$ entries after an eviction and
at most $n_0+1$ simultaneous reservations when the proposed and evicted
anchors are both included.
\end{lemma}

\begin{proof}
If $A_i$ precedes $A_j$, then $A_j$ lies between $A_i$ and $B$ on the
ancestry path.  Lemma~\ref{lem:separator}, applied to paths from $p\in B$
to $A_i$, gives $d(p,A_j)\leq d(p,A_i)$.  Thus a threshold protects at
most one entry.  There are at most $2k(k-1)$ protected entries in total.
At a full insertion there are $n_0$ non-root candidates, and
$n_0>2k(k-1)$, so at least one is eligible.  The size assertions follow
from performing exactly one eviction at such an insertion.
\end{proof}

\subsection{Realization and separation tests}

At an anchor $A$, use the ordinary perturbation of
Lemma~\ref{lem:ordinary-noise}:
\begin{equation}\label{eq:anchor-perturbation}
 h_A=d(\,\cdot\,,A),\qquad g_A=\xi_A h_A+v_A.
\end{equation}
All randomness at this launch is fresh.  Reserve the two particles
$f_A+g_A$ and $f_A-g_A$, where $f_A$ is the average of the selected pair
just before the launch.  Their reservation lasts through the eviction
bag, including that bag.  A rejected proposal therefore has a one-bag
reservation.  At an insertion bag, record the trace after the proposal's
launch and before releasing any reservation; the trace-preserving
operations of Proposition~\ref{prop:particle-realization} make this
convention consistent with both the incoming and outgoing lists.
Write $\rho>0$ and $c>0$ for the probability and amplitude constants,
respectively, in Lemma~\ref{lem:ordinary-noise}.

Set
\[
 q_0=n_0+1,\qquad N=2q_0+k!+2.
\]
This particle count covers the inclusive reservation bound in
Lemma~\ref{lem:cache-defined}.  Ordinary perturbations have Lipschitz
constant at most $3$ and absolute value at most $3\operatorname{diam}(V)$,
independently of $\eta$.  Proposition~\ref{prop:particle-realization}
therefore gives a cut pseudometric, averaged over all launch randomness, with
\begin{equation}\label{eq:anchor-upper}
 D_{\mathrm{anc}}\leq U d,\qquad U=3N(N-1).
\end{equation}
Every particle has Lipschitz constant at most
$L_{\mathrm{p}}=3(N-1)$.  The common sweeping interval can be chosen with
margin at least $2\operatorname{diam}(V)$ around all particle values.
In particular, $N$, $L_{\mathrm{p}}$, and $U$ depend only on $k$ and do
not depend on $J$ or $\eta$.

Write $b_k>0$ for the constant in
Proposition~\ref{prop:uncertainty}.

We will repeatedly use two consequences of the cut-measure lemmas.

\begin{samepage}
\begin{lemma}[Tests along a reservation]\label{lem:anchor-tests}
Let a path of bags run from a bag containing $y$ to a later bag containing
$z$.
\begin{enumerate}
\item If an anchor $A$ lies on this path and its reservation lasts from
$A$ through the later bag, then
\begin{equation}\label{eq:anchor-launch-test}
 D_{\mathrm{anc}}(y,z)\geq b_k d(z,A).
\end{equation}
\item Suppose a testing bag $P$ lies on the path and a reserved particle
$F$ is unchanged from $P$ through the later bag.  If, with probability at
least $\theta$, it satisfies
\[
 |F(z)-F(s)|\geq t\quad\text{for every }s\in P,
\]
where $0<t\leq\operatorname{diam}(V)$, then
\begin{equation}\label{eq:anchor-height-test}
 D_{\mathrm{anc}}(y,z)\geq b_k\theta t.
\end{equation}
In the second assertion the tested reservation and vertices are chosen
without consulting launch randomness.
\end{enumerate}
Both assertions allow the later bag to be the eviction bag.
\end{lemma}
\end{samepage}

\begin{proof}
For each realization of the launch randomness, a reserved pair supplies
a common submeasure throughout its reservation, including its endpoints.
Since $g_A$ vanishes on $A$ and $|g_A(z)|\geq h_A(z)/2$,
Lemma~\ref{lem:launch-gap} gives uncertainty at least $d(z,A)$.
Proposition~\ref{prop:uncertainty} gives the first assertion after
averaging.  For the second assertion, use the single particle's
symmetrized sweep.  The sweeping margin permits
Lemma~\ref{lem:height-gap} to give uncertainty at least $t$ on the stated
event.  Apply Proposition~\ref{prop:uncertainty} and average again.
These transport estimates hold pointwise in the prepared laws. To
establish the probability of a height gap in Section~\ref{sec:lower},
we will instead condition on the launch history and selected fresh
variables, leaving the later bag laws unconditioned.
\end{proof}

The preceding construction and tests hold for every $J>0$ and
$0<\eta<1/4$. We now fix the member of this family used in the rest
of the proof, choosing the parameters in the following order:
\begin{equation}\label{eq:anchor-parameters}
 K>2+\frac{2U+1}{b_k},\qquad
 J=10(K+1),\qquad H>10^5J,\qquad
 0<\eta<\frac{1}{10H}.
\end{equation}
The memory bound, and hence $U$, was fixed independently of $J$ and
$\eta$. All these parameters, including $\eta$, are now fixed for
every graph with bag size at most $k$. From this point onward,
$D_{\mathrm{anc}}$ denotes the pseudometric with these choices.
